% method_sim.tex
% Method: variance-corrected SIM composition.
% Simplified 2026-07-14 around a variance-budget explanation. The
% section now proceeds from the composition goal, to naive double
% counting, to the variance correction, the high-R^2 exception, and
% finally the evaluation protocol. Detailed verification remains in
% the Online Appendix.

The single-asset model describes each ticker separately. It does not
say how several tickers should move together. For the multi-asset
extension, we added one common source of movement: a market path. We used
the single-index model (SIM) of Sharpe~\cite{sharpe1963} as follows:
\begin{equation}
  g_i(t) \;=\; \alpha_i \;+\; \beta_i\, g_m(t) \;+\; \eps_i(t),
  \qquad t = 1,\dots,T.
  \label{eq:sim-composition}
\end{equation}
Here $\alpha_i$ is the intercept for asset $i$, $g_m$ is the market
growth-rate path, $\beta_i$ is the asset's market loading, and
$\eps_i$ is its asset-specific residual path. The symbol $\eps_i$
denotes a residual path;
it is unrelated to the jump probability $\epsilon$ used by the
single-asset model. We estimated the intercept and market loading
$(\alpha_i,\beta_i)$ from real data by
ordinary least squares (OLS). The remaining modeling choice is
$\eps_i$. A standard SIM draws independent Gaussian residuals with
the variance estimated by the real-data regression. This preserves
the fitted market relationship but discards the heavy tails and
volatility regimes produced by the single-asset generator. We called
this baseline \emph{Gaussian SIM}.

\paragraph{Why direct construction fails.}

A tempting alternative is to use a full synthetic return path
$\tg_i$ from the single-asset generator as the residual:
\begin{equation}\label{eq:naive-composition}
g_i = \alpha_i + \beta_i\, g_m + \tg_i.
\end{equation}
This raw substitution double counts location as well as variance. The
calibrated intercept already fixes the conditional location of asset $i$,
whereas the full-return draw generally has
$\E[\tg_i]\ne 0$. Equation~\eqref{eq:naive-composition} therefore has
conditional mean $\alpha_i+\E[\tg_i]+\beta_i g_m$, not
$\alpha_i+\beta_i g_m$.
To remove this location error, all full-return constructions in our
comparison first removed the realized path mean. We defined the
centered generator path as follows:
\begin{equation}
  \overline{\tg}_i \equiv \frac{1}{T}\sum_{t=1}^{T}\tg_i(t),
  \qquad
  \tg_i^{\mathrm c}(t) \equiv \tg_i(t)-\overline{\tg}_i.
  \label{eq:center-generator}
\end{equation}
Therefore, $T^{-1}\sum_t\tg_i^{\mathrm c}(t)=0$ exactly. Subtracting a
constant leaves both the sample variance and the sample covariance with
$g_m$ unchanged. We therefore used the centered direct construction
$g_i=\alpha_i+\beta_i g_m+\tg_i^{\mathrm c}$ as the \emph{naive}
variance comparator.

Centering removed the location error but did not correct the variance
error. The path variance
$\sigma_{\gen,i}^2=\Var(\tg_i)=\Var(\tg_i^{\mathrm c})$ is the
\emph{generator variance}. It already represents the asset's total
variance. Adding the market term therefore contributes variance a
second time. Independent sampling makes the residual--market covariance
zero in population and close to zero on a long realized path. Under
that working approximation, the naive composed variance is
\begin{equation}
  \Var(g_i)
  \;\approx\; \beta_i^2\sigma_m^2 + \sigma_{\gen,i}^2,
  \qquad \Cov(\tg_i^{\mathrm c},g_m)\approx 0,
  \label{eq:naive-inflation}
\end{equation}
where $\sigma_m^2=\Var(g_m)$ is the variance of the market path. The
excess is $\beta_i^2\sigma_m^2$, so the error is largest for
assets with high market exposure.

\paragraph{Variance-preserving correction.}

To prevent the market term from adding variance beyond the generator
target, we treated $\sigma_{\gen,i}^2$ as asset $i$'s total variance
budget. The market term uses part of that budget. We scaled the
single-asset draw so that it uses only the part that remains. We
defined $s_i$ as the nonnegative asset-specific residual scale, set
$\eps_i=s_i\tg_i^{\mathrm c}$, and required
$\Var(\beta_i g_m+s_i\tg_i^{\mathrm c})=\sigma_{\gen,i}^2$. Under the
same covariance approximation, solving for $s_i$ gives
\begin{equation}
  \boxed{\;s_i^2 \;=\; 1 \;-\; \frac{\beta_i^2\, \sigma_m^2}{\sigma_{\gen,i}^2}\;}
  \label{eq:scale}
\end{equation}
We defined the \emph{market loading ratio} as follows:
\begin{equation}
  \rho_i \;\equiv\; \frac{\beta_i^2\, \sigma_m^2}{\sigma_{\gen,i}^2},
  \label{eq:rho}
\end{equation}
Under the zero-covariance assumption, this ratio is the fraction of the
variance budget used by the market, and $s_i^2=1-\rho_i$. At
$\rho_i=1$, the unconstrained rule leaves no asset-specific variance;
for $\rho_i>1$, it has no real-valued solution. Either case can arise
when $\beta_i$ is large, the generator variance is low, or the
synthetic market is more volatile than the calibration market. To
retain a minimum asset-specific fraction $f$, we imposed the stricter
bound $\rho_i\leq 1-f$ and reduced the effective loading
$\beta_i^{\eff}$ whenever that bound was exceeded. The resulting rule
was:
\begin{equation}
  \beta_i^{\eff} \;=\;
  \begin{cases}
    \beta_i & \text{if } \rho_i \le \rho_i^{\max} \\[4pt]
    \sign(\beta_i)\,\sqrt{(1 - f)\,\dfrac{\sigma_{\gen,i}^2}{\sigma_m^2}} & \text{if } \rho_i > \rho_i^{\max}
  \end{cases}
  \qquad
  s_i^2 \;=\;
  \begin{cases}
    1 - \rho_i & \text{if } \rho_i \le \rho_i^{\max} \\[4pt]
    f          & \text{if } \rho_i > \rho_i^{\max}
  \end{cases}
  \label{eq:clip}
\end{equation}
The sign of the calibrated market loading $\beta_i$ is retained. We
used $f=0.10$, which guarantees on this ordinary-stock branch that at
least $10\%$ of the variance budget remains asset-specific
(Table~\ref{tab:hyperparams}).

\paragraph{High-$R^2$ index trackers.}

For an ordinary stock, our main target was total variance. For an
index-tracking fund, preserving its relationship with the market was
more important. The coefficient of determination $R^2$ measures the
fraction of the asset's variance explained by the market. For example,
SPY regressed on itself has $\beta=1$, $R^2=1$, and no residual
variance. The QQQ and SPYG funds also have high $R^2$ against SPY. Applying the
ordinary-stock rule to these trackers can preserve total variance
while changing their market relationship. We
therefore used a separate branch when the calibrated coefficient of
determination, $R^2_{i,\real}$, is at least the preservation threshold
$R^2_{\mathrm{preserve}}$. In our
universe, $0.80$ separated the index trackers from individual stocks
(Fig.~\ref{fig:r2_distribution}). On this branch, we computed the
target residual variance needed to recover the calibrated $R^2$ under
the same zero-covariance assumption:
\begin{equation}
  \boxed{\;\sigma_{\eps,\mathrm{target}}^2 \;=\; \beta_i^2\, \sigma_m^2 \cdot
    \frac{1 - R^2_{i,\real}}{R^2_{i,\real}}\;}
  \label{eq:r2-target}
\end{equation}
We rescaled the generator draw to this target as follows:
\begin{equation}
  \eps_i \;=\; \left(\sqrt{\frac{\sigma_{\eps,\mathrm{target}}^2}{\sigma_{\gen,i}^2}}\,\right)\tg_i^{\mathrm c}.
  \label{eq:r2-rescale}
\end{equation}
When $R^2_{i,\real}=1$, the target residual variance is zero, as it
should be for SPY against itself. If the scale factor in
Eq.~\eqref{eq:r2-rescale} exceeds one, the procedure stretches the
generator draw and can weaken the original marginal fit.

\paragraph{Complete composition rule.}

Together, these steps targeted the calibrated $R^2$ for a high-$R^2$
tracker and the generator variance for other assets. For the latter,
the method clipped beta only when the market term would leave too
little asset-specific variance, then added the scaled asset draw to
the market term (Algorithm~\ref{alg:hybrid}). A residual-copula rank
reorder, if used, is applied jointly after scaling and before final
composition; it preserves each residual path's empirical variance.
Centering kept the path mean controlled by the intercept and market
term. Under the zero-covariance assumption, an OLS regression on a
long composed path therefore converges to the effective market loading
and its corresponding intercept. The ordinary-stock branch targeted
total generator variance; the tracker branch targeted residual
variance and thus the total variance implied by $(\beta_i,
R^2_{i,\real})$. These are asymptotic targets, not exact finite-sample
guarantees or claims about the complete return distribution. Online
Appendix~\ref{sec:supp-sim-derivations} gives the algebra. The
construction adds only same-day market dependence; it does not
reproduce lead--lag effects or market-driven volatility clustering.

\begin{algorithm}[t]
\caption{Hybrid SIM composition with variance correction.}
\label{alg:hybrid}
\small
\begin{algorithmic}[1]
  \REQUIRE Nonconstant $g_m,\tg_i\in\R^T$, $i=1,\dots,N$; each
    $\tg_i$ is independent of $g_m$.
  \REQUIRE Parameters $(\alpha_i,\beta_i,R^2_{i,\real})$; floor
    $f\in(0,1)$; threshold $R^2_{\mathrm{preserve}}\in(0,1]$.
  \ENSURE Composed paths $g_i\in\R^T$ and branch flags $b_i$.
  \STATE $\sigma_m^2 \gets \Var(g_m)$.
  \FOR{each asset $i = 1,\dots,N$}
    \STATE $\sigma_{\gen,i}^2 \gets \Var(\tg_i)$;
      $\tg_i^{\mathrm c} \gets \tg_i-\overline{\tg}_i$.
    \IF{$R^2_{i,\real} \ge R^2_{\mathrm{preserve}}$}
      \STATE $v_i \gets \beta_i^2\sigma_m^2(1-R^2_{i,\real})/R^2_{i,\real}$.
      \STATE $s_i \gets \sqrt{v_i/\sigma_{\gen,i}^2}$;
        $\beta_i^{\eff} \gets \beta_i$; $b_i\gets\textsc{r2-preserve}$.
    \ELSE
      \STATE $\rho_i \gets \beta_i^2 \sigma_m^2 / \sigma_{\gen,i}^2$
      \IF{$\rho_i \le 1 - f$}
        \STATE $s_i \gets \sqrt{1-\rho_i}$;
          $\beta_i^{\eff}\gets\beta_i$; $b_i\gets\textsc{variance}$.
      \ELSE
        \STATE $s_i\gets\sqrt{f}$;
          $\beta_i^{\eff}\gets\sign(\beta_i)
          \sqrt{(1-f)\sigma_{\gen,i}^2/\sigma_m^2}$.
        \STATE $b_i\gets\textsc{clipped}$.
      \ENDIF
    \ENDIF
    \STATE $\eps_i \gets s_i\tg_i^{\mathrm c}$.
  \ENDFOR
  \STATE \textit{Optional:} jointly rank-reorder $\{\eps_i\}_{i=1}^{N}$.
  \FOR{each asset $i = 1,\dots,N$}
    \STATE $g_i \gets \alpha_i + \beta_i^{\eff}\, g_m + \eps_i$
      \hfill \COMMENT{Eq.~\eqref{eq:sim-composition}}
  \ENDFOR
\end{algorithmic}
\end{algorithm}

% --- Table: Hyperparameters ---
\begin{table}[tbp]
\centering
\caption{\textbf{HMM-WJ and SIM composition hyperparameters.} Values
    used for the main SPY analysis, with the sensitivity-sweep or
    grid-search span and the location of the corresponding check
    where one exists. The $423$-ticker SIM composition universe holds
    $\epsilon = 0$ instead of the SPY value; the cross-asset
    validation used per-ticker
    $(\epsilon^{*}, \lambda^{*})$ for the NVDA, JNJ, and JPM
    models (Table~\ref{tab:cross_asset}).}
\label{tab:hyperparams}
\small
\begin{tabular}{@{}llll@{}}
\toprule
\textbf{Symbol} & \textbf{Description} & \textbf{SPY value} & \textbf{Explored range} \\ \midrule
\multicolumn{4}{l}{\textit{State partition and emission}} \\[2pt]
$N$            & Laplace-quantile states                      & $100$     & $\{30,\dots,350\}$, Online Appendix~\ref{sec:supp-sensitivity} \\
$\nu$          & Student-$t$ emission degrees of freedom      & $5$       & $\{3,\dots,30,\infty\}$ \\[4pt]
\multicolumn{4}{l}{\textit{Jump-duration override}} \\[2pt]
$\epsilon$     & jump probability per step                    & $10^{-4}$ & $[10^{-4}, 2.5\times10^{-2}]$, Eq.~\eqref{eq:grid_search} \\
$\lambda$      & mean jump-episode duration (steps)           & $90$      & $[10, 160]$, Eq.~\eqref{eq:grid_search} \\
$p_{\rm neg}$  & probability a jump targets $\mathcal{S}_{-}$ & $0.52$    & fixed \\
$N_{\rm tail}$ & quantile states per tail set                 & $5$       & fixed (at $N = 100$) \\
$w_\kappa$     & kurtosis-penalty weight, Eq.~\eqref{eq:grid_search} & $0.20$ & fixed \\
$L$            & ACF lag window, Eq.~\eqref{eq:grid_search}   & $25$ days & fixed \\[4pt]
\multicolumn{4}{l}{\textit{SIM composition}} \\[2pt]
$f$                        & idiosyncratic variance floor & $0.10$ & $\{0.05,\dots,0.30\}$ \\
$R^2_{\mathrm{preserve}}$  & branch-selection threshold   & $0.80$ & $\{0.70,\dots,0.90\}$ \\
\bottomrule
\end{tabular}
\end{table}
\FloatBarrier

\paragraph{Evaluation protocol.}

To test whether the composition method scaled across a broad equity
universe, we used a complete-case sample of $424$ US equities and
exchange-traded funds from Polygon.io, each with $2{,}767$ daily
closing prices from January 2014 through December $2024$. This
criterion yielded $2{,}766$ growth rates per ticker and conditioned the
analysis on assets that survived with full training-window coverage.
We treated SPY as the market and composed the remaining $423$ assets,
using the SPY state and emission specification for each
(Table~\ref{tab:hyperparams}; $r_f=0$). To isolate multi-asset
composition from per-ticker jump calibration, we set the jump
probability to zero. Regressing each real asset on SPY then estimated
$\alpha_i$, $\beta_i$, $R^2_{i,\real}$, and residual variance. The
resulting $\beta$ range of $0.03$--$1.79$ and $R^2_{i,\real}$ range of
$0.001$--$0.933$ exposed the methods to widely differing market
dependence (Fig.~\ref{fig:r2_distribution}).

To separate the effect of variance correction from the choice of
asset-specific generator, we compared six construction methods. The naive
method used the full JumpHMM draw without correction. The Gaussian
SIM method used independent Gaussian residuals with variance estimated
by OLS. The hybrid method applied the variance-corrected composition
described above (Algorithm~\ref{alg:hybrid}). The
JumpHMM-on-residuals method fitted a new JumpHMM to the OLS residual
series after converting the residuals to a synthetic price path by
cumulative compounding. The block-bootstrap method sampled the OLS
residual series with the Politis--Romano stationary
bootstrap~\cite{politisRomano1994} and a mean block length of
$\sqrt{T}\approx 50$. The GARCH(1,1)-$t$ method fitted a
conditional-variance model to each residual series~\cite{bollerslev1986}.
We excluded thirty non-stationary fits from the GARCH results only.
For each ticker, the naive and hybrid methods received the same
JumpHMM draw. This pairing isolated the effect of the variance
correction. The other four methods generated their own residual draws.
We did not use the optional cross-sectional rank reorder because the
experiment tested each ticker's marginal construction separately. This
design made the naive--hybrid pairing a direct test of the correction
and the remaining methods a test of whether comparable performance
required a separately fitted residual generator.

To determine whether each method retained both the calibrated market
relationship and the asset generator it reused, we evaluated every
composed series along three dimensions. First, a new OLS regression
tested recovery of the calibrated market loading and $R^2$ through the
absolute errors $|\hat\beta-\beta|$ and $|\hat R^2-R^2_{\real}|$,
where lower values are better. Second, we compared the synthetic and
real return distributions using
Kolmogorov--Smirnov (KS) and Anderson--Darling (AD) pass rates,
Wasserstein-1 distance, excess kurtosis, and the Hill upper-tail index.
Higher KS and AD pass rates and lower Wasserstein distance indicate a
closer match; kurtosis and the Hill index describe tail behavior.
Third, a composed-to-generator variance ratio tested preservation of
generator variance, with one as the ideal value.
Online Appendix~\ref{sec:supp-validation} gives the metric definitions.
We ran $100$ replications per ticker with seed $1234$, producing
$42{,}300$ series per method. Pass rates used the significance level
$\alpha=0.05$. Together, these measures formed the in-sample scorecard
and the basis for the frozen-parameter holdout comparison.

To test whether performance on that scorecard generalized without
refitting or look-ahead bias, we froze all marginal fits, residual
models, and SIM parameters
estimated from $2014$--$2024$ and evaluated the methods against the
$249$ growth rates observed in $2025$. Of the $423$ non-market
training assets, $416$ had complete holdout histories; GARCH results
covered $386$ because its thirty non-stationary training fits remained
excluded. For each replication, we generated a $249$-day SPY path
from the training-period SPY marginal and used that common simulated
market path in all asset compositions. Observed $2025$ SPY and asset
returns were used only for scoring. We repeated the KS, AD,
Wasserstein-1, variance, kurtosis, and one-day left-tail VaR evaluations with
$100$ paths per ticker. Because goodness-of-fit tests have less power
on $249$ observations than on $2{,}766$, we also compared each
synthetic path with a randomly selected contiguous $249$-day block
from the training history. We first averaged replication-level
results within ticker and then summarized across tickers so each
asset received equal weight. This design made the $2025$ comparison
strictly out of sample while controlling for the lower power of its
shorter histories.
